% Architecture diagram for clio-agentic-search
% Usage: % Architecture diagram for clio-agentic-search
% Usage: % Architecture diagram for clio-agentic-search
% Usage: % Architecture diagram for clio-agentic-search
% Usage: \input{figures/architecture.tex}
\begin{figure}[!t]
\centering
\begin{tikzpicture}[
  % --- global style ---
  >=Stealth,
  node distance=0.45cm and 0.25cm,
  every node/.style={font=\scriptsize},
  box/.style={
    draw, rounded corners=2pt, minimum height=0.55cm,
    text centered, inner sep=2pt, fill=white,
    minimum width=2.0cm
  },
  scibox/.style={
    box, fill=black!8
  },
  connbox/.style={
    draw, rounded corners=2pt, minimum height=0.45cm,
    text centered, inner sep=2pt, fill=white,
    minimum width=1.5cm, font=\tiny
  },
  sciconn/.style={
    connbox, fill=black!8
  },
  arr/.style={->, thick},
  darr/.style={->, densely dashed, thick},
]

% ====== Main pipeline (centred) ======
\node[box, minimum width=2.8cm] (query) {Query Interface};
\node[below=of query, font=\tiny] (querysub) {FastAPI\,/\,CLI};

\node[box, minimum width=2.8cm, below=0.55cm of query] (ns) {Namespace Registry};

\node[box, minimum width=2.8cm, below=of ns] (rc) {Retrieval Coordinator};

% --- four parallel branches ---
\node[box, below=0.7cm of rc, xshift=-1.55cm, minimum width=1.15cm]
  (lex) {Lexical};
\node[below=0pt of lex, font=\tiny] (lexsub) {BM25};

\node[box, below=0.7cm of rc, xshift=-0.5cm, minimum width=1.15cm]
  (vec) {Vector};
\node[below=0pt of vec, font=\tiny] (vecsub) {Dense};

\node[box, below=0.7cm of rc, xshift=0.55cm, minimum width=1.15cm]
  (gra) {Graph};
\node[below=0pt of gra, font=\tiny] (grasub) {BFS};

\node[scibox, below=0.7cm of rc, xshift=1.6cm, minimum width=1.15cm]
  (sci) {Scientific};
\node[below=0pt of sci, font=\tiny] (scisub) {SI\,+\,Formula};

% --- merge + rerank ---
\node[box, minimum width=2.8cm, below=1.15cm of rc, yshift=-0.95cm]
  (merge) {Merge + Rerank};

% --- citations ---
\node[box, minimum width=2.8cm, below=of merge]
  (cite) {Citations + Traces};

% ====== Arrows: main pipeline ======
\draw[arr] (query) -- (ns);
\draw[arr] (ns) -- (rc);

% fan-out
\draw[arr] (rc.south) -- ++(0,-0.25) -| (lex.north);
\draw[arr] (rc.south) -- ++(0,-0.25) -| (vec.north);
\draw[arr] (rc.south) -- ++(0,-0.25) -| (gra.north);
\draw[arr] (rc.south) -- ++(0,-0.25) -| (sci.north);

% fan-in
\draw[arr] (lex.south) ++(0,-0.15) -- ++(0,-0.15) -| (merge.north -| lex.south) -- (merge.north);
\draw[arr] (vec.south) ++(0,-0.15) -- ++(0,-0.15) -| (merge.north -| vec.south) -- (merge.north);
\draw[arr] (gra.south) ++(0,-0.15) -- ++(0,-0.15) -| (merge.north -| gra.south) -- (merge.north);
\draw[arr] (sci.south) ++(0,-0.15) -- ++(0,-0.15) -| (merge.north -| sci.south) -- (merge.north);

\draw[arr] (merge) -- (cite);

% ====== Right side: connectors ======
\node[connbox, right=1.1cm of ns, yshift=0.85cm]  (fs)   {Filesystem};
\node[connbox, below=0.10cm of fs]                 (s3)   {S3 Object Store};
\node[connbox, below=0.10cm of s3]                 (qd)   {Qdrant};
\node[connbox, below=0.10cm of qd]                 (neo)  {Neo4j};
\node[sciconn, below=0.10cm of neo]                (hdf)  {HDF5};
\node[sciconn, below=0.10cm of hdf]                (nc)   {NetCDF};
\node[sciconn, below=0.10cm of nc]                 (iow)  {IOWarp CTE};
\node[connbox, below=0.10cm of iow]                (ndp)  {NDP (MCP)};

% label
\node[above=0.08cm of fs, font=\tiny\bfseries] {Connectors (9)};

% arrows from connectors to namespace registry
\draw[arr] (fs.west)  -- (ns.east |- fs.west);
\draw[arr] (s3.west)  -- (ns.east |- s3.west);
\draw[arr] (qd.west)  -| ([xshift=0.15cm]ns.east) -- (ns.east);
\draw[arr] (neo.west) -| ([xshift=0.15cm]ns.east) -- (ns.east);
\draw[arr] (hdf.west) -| ([xshift=0.15cm]ns.east) -- (ns.east);
\draw[arr] (nc.west)  -| ([xshift=0.15cm]ns.east) -- (ns.east);
\draw[arr] (iow.west) -| ([xshift=0.15cm]ns.east) -- (ns.east);
\draw[arr] (ndp.west) -| ([xshift=0.15cm]ns.east) -- (ns.east);

% ====== Left side: agentic loop ======
\node[box, left=1.0cm of rc, minimum width=1.6cm, fill=white]
  (llm) {\parbox{1.4cm}{\centering LLM Query\\[-1pt] Rewriter}};

% agentic loop arrows
\draw[darr] (cite.west) -- ++(-0.35,0) |- (llm.south);
\draw[darr] (llm.east) -- (rc.west);

% loop label
\node[left=0.04cm of llm, font=\tiny\bfseries, rotate=90, anchor=south]
  {Agentic Loop};

% background highlight for agentic loop
\begin{pgfonlayer}{background}
  \node[draw=black!40, dashed, rounded corners=4pt, fill=black!3,
        inner sep=7pt,
        fit=(llm)(rc)(merge)(cite)] {};
\end{pgfonlayer}

\end{tikzpicture}
\caption{System architecture of \textsc{clio-agentic-search}. Science-aware
components (shaded) execute alongside standard retrieval branches. The agentic
loop iteratively refines queries via LLM rewriting.}
\label{fig:architecture}
\end{figure}

\begin{figure}[!t]
\centering
\begin{tikzpicture}[
  % --- global style ---
  >=Stealth,
  node distance=0.45cm and 0.25cm,
  every node/.style={font=\scriptsize},
  box/.style={
    draw, rounded corners=2pt, minimum height=0.55cm,
    text centered, inner sep=2pt, fill=white,
    minimum width=2.0cm
  },
  scibox/.style={
    box, fill=black!8
  },
  connbox/.style={
    draw, rounded corners=2pt, minimum height=0.45cm,
    text centered, inner sep=2pt, fill=white,
    minimum width=1.5cm, font=\tiny
  },
  sciconn/.style={
    connbox, fill=black!8
  },
  arr/.style={->, thick},
  darr/.style={->, densely dashed, thick},
]

% ====== Main pipeline (centred) ======
\node[box, minimum width=2.8cm] (query) {Query Interface};
\node[below=of query, font=\tiny] (querysub) {FastAPI\,/\,CLI};

\node[box, minimum width=2.8cm, below=0.55cm of query] (ns) {Namespace Registry};

\node[box, minimum width=2.8cm, below=of ns] (rc) {Retrieval Coordinator};

% --- four parallel branches ---
\node[box, below=0.7cm of rc, xshift=-1.55cm, minimum width=1.15cm]
  (lex) {Lexical};
\node[below=0pt of lex, font=\tiny] (lexsub) {BM25};

\node[box, below=0.7cm of rc, xshift=-0.5cm, minimum width=1.15cm]
  (vec) {Vector};
\node[below=0pt of vec, font=\tiny] (vecsub) {Dense};

\node[box, below=0.7cm of rc, xshift=0.55cm, minimum width=1.15cm]
  (gra) {Graph};
\node[below=0pt of gra, font=\tiny] (grasub) {BFS};

\node[scibox, below=0.7cm of rc, xshift=1.6cm, minimum width=1.15cm]
  (sci) {Scientific};
\node[below=0pt of sci, font=\tiny] (scisub) {SI\,+\,Formula};

% --- merge + rerank ---
\node[box, minimum width=2.8cm, below=1.15cm of rc, yshift=-0.95cm]
  (merge) {Merge + Rerank};

% --- citations ---
\node[box, minimum width=2.8cm, below=of merge]
  (cite) {Citations + Traces};

% ====== Arrows: main pipeline ======
\draw[arr] (query) -- (ns);
\draw[arr] (ns) -- (rc);

% fan-out
\draw[arr] (rc.south) -- ++(0,-0.25) -| (lex.north);
\draw[arr] (rc.south) -- ++(0,-0.25) -| (vec.north);
\draw[arr] (rc.south) -- ++(0,-0.25) -| (gra.north);
\draw[arr] (rc.south) -- ++(0,-0.25) -| (sci.north);

% fan-in
\draw[arr] (lex.south) ++(0,-0.15) -- ++(0,-0.15) -| (merge.north -| lex.south) -- (merge.north);
\draw[arr] (vec.south) ++(0,-0.15) -- ++(0,-0.15) -| (merge.north -| vec.south) -- (merge.north);
\draw[arr] (gra.south) ++(0,-0.15) -- ++(0,-0.15) -| (merge.north -| gra.south) -- (merge.north);
\draw[arr] (sci.south) ++(0,-0.15) -- ++(0,-0.15) -| (merge.north -| sci.south) -- (merge.north);

\draw[arr] (merge) -- (cite);

% ====== Right side: connectors ======
\node[connbox, right=1.1cm of ns, yshift=0.85cm]  (fs)   {Filesystem};
\node[connbox, below=0.10cm of fs]                 (s3)   {S3 Object Store};
\node[connbox, below=0.10cm of s3]                 (qd)   {Qdrant};
\node[connbox, below=0.10cm of qd]                 (neo)  {Neo4j};
\node[sciconn, below=0.10cm of neo]                (hdf)  {HDF5};
\node[sciconn, below=0.10cm of hdf]                (nc)   {NetCDF};
\node[sciconn, below=0.10cm of nc]                 (iow)  {IOWarp CTE};
\node[connbox, below=0.10cm of iow]                (ndp)  {NDP (MCP)};

% label
\node[above=0.08cm of fs, font=\tiny\bfseries] {Connectors (9)};

% arrows from connectors to namespace registry
\draw[arr] (fs.west)  -- (ns.east |- fs.west);
\draw[arr] (s3.west)  -- (ns.east |- s3.west);
\draw[arr] (qd.west)  -| ([xshift=0.15cm]ns.east) -- (ns.east);
\draw[arr] (neo.west) -| ([xshift=0.15cm]ns.east) -- (ns.east);
\draw[arr] (hdf.west) -| ([xshift=0.15cm]ns.east) -- (ns.east);
\draw[arr] (nc.west)  -| ([xshift=0.15cm]ns.east) -- (ns.east);
\draw[arr] (iow.west) -| ([xshift=0.15cm]ns.east) -- (ns.east);
\draw[arr] (ndp.west) -| ([xshift=0.15cm]ns.east) -- (ns.east);

% ====== Left side: agentic loop ======
\node[box, left=1.0cm of rc, minimum width=1.6cm, fill=white]
  (llm) {\parbox{1.4cm}{\centering LLM Query\\[-1pt] Rewriter}};

% agentic loop arrows
\draw[darr] (cite.west) -- ++(-0.35,0) |- (llm.south);
\draw[darr] (llm.east) -- (rc.west);

% loop label
\node[left=0.04cm of llm, font=\tiny\bfseries, rotate=90, anchor=south]
  {Agentic Loop};

% background highlight for agentic loop
\begin{pgfonlayer}{background}
  \node[draw=black!40, dashed, rounded corners=4pt, fill=black!3,
        inner sep=7pt,
        fit=(llm)(rc)(merge)(cite)] {};
\end{pgfonlayer}

\end{tikzpicture}
\caption{System architecture of \textsc{clio-agentic-search}. Science-aware
components (shaded) execute alongside standard retrieval branches. The agentic
loop iteratively refines queries via LLM rewriting.}
\label{fig:architecture}
\end{figure}

\begin{figure}[!t]
\centering
\begin{tikzpicture}[
  % --- global style ---
  >=Stealth,
  node distance=0.45cm and 0.25cm,
  every node/.style={font=\scriptsize},
  box/.style={
    draw, rounded corners=2pt, minimum height=0.55cm,
    text centered, inner sep=2pt, fill=white,
    minimum width=2.0cm
  },
  scibox/.style={
    box, fill=black!8
  },
  connbox/.style={
    draw, rounded corners=2pt, minimum height=0.45cm,
    text centered, inner sep=2pt, fill=white,
    minimum width=1.5cm, font=\tiny
  },
  sciconn/.style={
    connbox, fill=black!8
  },
  arr/.style={->, thick},
  darr/.style={->, densely dashed, thick},
]

% ====== Main pipeline (centred) ======
\node[box, minimum width=2.8cm] (query) {Query Interface};
\node[below=of query, font=\tiny] (querysub) {FastAPI\,/\,CLI};

\node[box, minimum width=2.8cm, below=0.55cm of query] (ns) {Namespace Registry};

\node[box, minimum width=2.8cm, below=of ns] (rc) {Retrieval Coordinator};

% --- four parallel branches ---
\node[box, below=0.7cm of rc, xshift=-1.55cm, minimum width=1.15cm]
  (lex) {Lexical};
\node[below=0pt of lex, font=\tiny] (lexsub) {BM25};

\node[box, below=0.7cm of rc, xshift=-0.5cm, minimum width=1.15cm]
  (vec) {Vector};
\node[below=0pt of vec, font=\tiny] (vecsub) {Dense};

\node[box, below=0.7cm of rc, xshift=0.55cm, minimum width=1.15cm]
  (gra) {Graph};
\node[below=0pt of gra, font=\tiny] (grasub) {BFS};

\node[scibox, below=0.7cm of rc, xshift=1.6cm, minimum width=1.15cm]
  (sci) {Scientific};
\node[below=0pt of sci, font=\tiny] (scisub) {SI\,+\,Formula};

% --- merge + rerank ---
\node[box, minimum width=2.8cm, below=1.15cm of rc, yshift=-0.95cm]
  (merge) {Merge + Rerank};

% --- citations ---
\node[box, minimum width=2.8cm, below=of merge]
  (cite) {Citations + Traces};

% ====== Arrows: main pipeline ======
\draw[arr] (query) -- (ns);
\draw[arr] (ns) -- (rc);

% fan-out
\draw[arr] (rc.south) -- ++(0,-0.25) -| (lex.north);
\draw[arr] (rc.south) -- ++(0,-0.25) -| (vec.north);
\draw[arr] (rc.south) -- ++(0,-0.25) -| (gra.north);
\draw[arr] (rc.south) -- ++(0,-0.25) -| (sci.north);

% fan-in
\draw[arr] (lex.south) ++(0,-0.15) -- ++(0,-0.15) -| (merge.north -| lex.south) -- (merge.north);
\draw[arr] (vec.south) ++(0,-0.15) -- ++(0,-0.15) -| (merge.north -| vec.south) -- (merge.north);
\draw[arr] (gra.south) ++(0,-0.15) -- ++(0,-0.15) -| (merge.north -| gra.south) -- (merge.north);
\draw[arr] (sci.south) ++(0,-0.15) -- ++(0,-0.15) -| (merge.north -| sci.south) -- (merge.north);

\draw[arr] (merge) -- (cite);

% ====== Right side: connectors ======
\node[connbox, right=1.1cm of ns, yshift=0.85cm]  (fs)   {Filesystem};
\node[connbox, below=0.10cm of fs]                 (s3)   {S3 Object Store};
\node[connbox, below=0.10cm of s3]                 (qd)   {Qdrant};
\node[connbox, below=0.10cm of qd]                 (neo)  {Neo4j};
\node[sciconn, below=0.10cm of neo]                (hdf)  {HDF5};
\node[sciconn, below=0.10cm of hdf]                (nc)   {NetCDF};
\node[sciconn, below=0.10cm of nc]                 (iow)  {IOWarp CTE};
\node[connbox, below=0.10cm of iow]                (ndp)  {NDP (MCP)};

% label
\node[above=0.08cm of fs, font=\tiny\bfseries] {Connectors (9)};

% arrows from connectors to namespace registry
\draw[arr] (fs.west)  -- (ns.east |- fs.west);
\draw[arr] (s3.west)  -- (ns.east |- s3.west);
\draw[arr] (qd.west)  -| ([xshift=0.15cm]ns.east) -- (ns.east);
\draw[arr] (neo.west) -| ([xshift=0.15cm]ns.east) -- (ns.east);
\draw[arr] (hdf.west) -| ([xshift=0.15cm]ns.east) -- (ns.east);
\draw[arr] (nc.west)  -| ([xshift=0.15cm]ns.east) -- (ns.east);
\draw[arr] (iow.west) -| ([xshift=0.15cm]ns.east) -- (ns.east);
\draw[arr] (ndp.west) -| ([xshift=0.15cm]ns.east) -- (ns.east);

% ====== Left side: agentic loop ======
\node[box, left=1.0cm of rc, minimum width=1.6cm, fill=white]
  (llm) {\parbox{1.4cm}{\centering LLM Query\\[-1pt] Rewriter}};

% agentic loop arrows
\draw[darr] (cite.west) -- ++(-0.35,0) |- (llm.south);
\draw[darr] (llm.east) -- (rc.west);

% loop label
\node[left=0.04cm of llm, font=\tiny\bfseries, rotate=90, anchor=south]
  {Agentic Loop};

% background highlight for agentic loop
\begin{pgfonlayer}{background}
  \node[draw=black!40, dashed, rounded corners=4pt, fill=black!3,
        inner sep=7pt,
        fit=(llm)(rc)(merge)(cite)] {};
\end{pgfonlayer}

\end{tikzpicture}
\caption{System architecture of \textsc{clio-agentic-search}. Science-aware
components (shaded) execute alongside standard retrieval branches. The agentic
loop iteratively refines queries via LLM rewriting.}
\label{fig:architecture}
\end{figure}

\begin{figure}[!t]
\centering
\begin{tikzpicture}[
  % --- global style ---
  >=Stealth,
  node distance=0.45cm and 0.25cm,
  every node/.style={font=\scriptsize},
  box/.style={
    draw, rounded corners=2pt, minimum height=0.55cm,
    text centered, inner sep=2pt, fill=white,
    minimum width=2.0cm
  },
  scibox/.style={
    box, fill=black!8
  },
  connbox/.style={
    draw, rounded corners=2pt, minimum height=0.45cm,
    text centered, inner sep=2pt, fill=white,
    minimum width=1.5cm, font=\tiny
  },
  sciconn/.style={
    connbox, fill=black!8
  },
  arr/.style={->, thick},
  darr/.style={->, densely dashed, thick},
]

% ====== Main pipeline (centred) ======
\node[box, minimum width=2.8cm] (query) {Query Interface};
\node[below=of query, font=\tiny] (querysub) {FastAPI\,/\,CLI};

\node[box, minimum width=2.8cm, below=0.55cm of query] (ns) {Namespace Registry};

\node[box, minimum width=2.8cm, below=of ns] (rc) {Retrieval Coordinator};

% --- four parallel branches ---
\node[box, below=0.7cm of rc, xshift=-1.55cm, minimum width=1.15cm]
  (lex) {Lexical};
\node[below=0pt of lex, font=\tiny] (lexsub) {BM25};

\node[box, below=0.7cm of rc, xshift=-0.5cm, minimum width=1.15cm]
  (vec) {Vector};
\node[below=0pt of vec, font=\tiny] (vecsub) {Dense};

\node[box, below=0.7cm of rc, xshift=0.55cm, minimum width=1.15cm]
  (gra) {Graph};
\node[below=0pt of gra, font=\tiny] (grasub) {BFS};

\node[scibox, below=0.7cm of rc, xshift=1.6cm, minimum width=1.15cm]
  (sci) {Scientific};
\node[below=0pt of sci, font=\tiny] (scisub) {SI\,+\,Formula};

% --- merge + rerank ---
\node[box, minimum width=2.8cm, below=1.15cm of rc, yshift=-0.95cm]
  (merge) {Merge + Rerank};

% --- citations ---
\node[box, minimum width=2.8cm, below=of merge]
  (cite) {Citations + Traces};

% ====== Arrows: main pipeline ======
\draw[arr] (query) -- (ns);
\draw[arr] (ns) -- (rc);

% fan-out
\draw[arr] (rc.south) -- ++(0,-0.25) -| (lex.north);
\draw[arr] (rc.south) -- ++(0,-0.25) -| (vec.north);
\draw[arr] (rc.south) -- ++(0,-0.25) -| (gra.north);
\draw[arr] (rc.south) -- ++(0,-0.25) -| (sci.north);

% fan-in
\draw[arr] (lex.south) ++(0,-0.15) -- ++(0,-0.15) -| (merge.north -| lex.south) -- (merge.north);
\draw[arr] (vec.south) ++(0,-0.15) -- ++(0,-0.15) -| (merge.north -| vec.south) -- (merge.north);
\draw[arr] (gra.south) ++(0,-0.15) -- ++(0,-0.15) -| (merge.north -| gra.south) -- (merge.north);
\draw[arr] (sci.south) ++(0,-0.15) -- ++(0,-0.15) -| (merge.north -| sci.south) -- (merge.north);

\draw[arr] (merge) -- (cite);

% ====== Right side: connectors ======
\node[connbox, right=1.1cm of ns, yshift=0.85cm]  (fs)   {Filesystem};
\node[connbox, below=0.10cm of fs]                 (s3)   {S3 Object Store};
\node[connbox, below=0.10cm of s3]                 (qd)   {Qdrant};
\node[connbox, below=0.10cm of qd]                 (neo)  {Neo4j};
\node[sciconn, below=0.10cm of neo]                (hdf)  {HDF5};
\node[sciconn, below=0.10cm of hdf]                (nc)   {NetCDF};
\node[sciconn, below=0.10cm of nc]                 (iow)  {IOWarp CTE};
\node[connbox, below=0.10cm of iow]                (ndp)  {NDP (MCP)};

% label
\node[above=0.08cm of fs, font=\tiny\bfseries] {Connectors (9)};

% arrows from connectors to namespace registry
\draw[arr] (fs.west)  -- (ns.east |- fs.west);
\draw[arr] (s3.west)  -- (ns.east |- s3.west);
\draw[arr] (qd.west)  -| ([xshift=0.15cm]ns.east) -- (ns.east);
\draw[arr] (neo.west) -| ([xshift=0.15cm]ns.east) -- (ns.east);
\draw[arr] (hdf.west) -| ([xshift=0.15cm]ns.east) -- (ns.east);
\draw[arr] (nc.west)  -| ([xshift=0.15cm]ns.east) -- (ns.east);
\draw[arr] (iow.west) -| ([xshift=0.15cm]ns.east) -- (ns.east);
\draw[arr] (ndp.west) -| ([xshift=0.15cm]ns.east) -- (ns.east);

% ====== Left side: agentic loop ======
\node[box, left=1.0cm of rc, minimum width=1.6cm, fill=white]
  (llm) {\parbox{1.4cm}{\centering LLM Query\\[-1pt] Rewriter}};

% agentic loop arrows
\draw[darr] (cite.west) -- ++(-0.35,0) |- (llm.south);
\draw[darr] (llm.east) -- (rc.west);

% loop label
\node[left=0.04cm of llm, font=\tiny\bfseries, rotate=90, anchor=south]
  {Agentic Loop};

% background highlight for agentic loop
\begin{pgfonlayer}{background}
  \node[draw=black!40, dashed, rounded corners=4pt, fill=black!3,
        inner sep=7pt,
        fit=(llm)(rc)(merge)(cite)] {};
\end{pgfonlayer}

\end{tikzpicture}
\caption{System architecture of \textsc{clio-agentic-search}. Science-aware
components (shaded) execute alongside standard retrieval branches. The agentic
loop iteratively refines queries via LLM rewriting.}
\label{fig:architecture}
\end{figure}
